% Lösungen Einheit 5 von 8 – Mammutbäume (zusammenbau v0.1)
\einheitenkopf{Einheit 5 von 8 · Mammutbäume}

\erg{23}{a) mehrere Pfade mit denselben Faktoren (3 Stück) \quad b) $3 \cdot 0{,}2 \cdot 0{,}8^2 = 0{,}384$ \quad c) $6 \cdot 0{,}5 \cdot 0{,}3 \cdot 0{,}2 = 0{,}18$ \quad d) $0{,}3^3 + 6 \cdot 0{,}5 \cdot 0{,}3 \cdot 0{,}2 = 0{,}207$ \quad e) $4 \cdot 0{,}6^3 \cdot 0{,}4^3 \approx 0{,}0553$ \quad f) $\frac{108}{210} + \frac{24}{210} = \frac{22}{35}$ \quad g) „4 über 2“ · „6 über 1“ : „10 über 3“ $= \frac{36}{120} = \frac{3}{10}$ \quad h) $n \cdot \frac{1}{6} \cdot \left(\frac{1}{2}\right)^{n-1}$ (genau eine 3, sonst 1, n Positionen); für $n = 4$: $4 \cdot \frac{1}{6} \cdot \frac{1}{8} = \frac{1}{12}$}
\erg{24}{a) Max vergisst die drei Reihenfolgen: $3 \cdot \frac{5}{216} = \frac{5}{72}$ \quad b) Die Faktoren sind dieselben, nur an anderer Stelle; Zähler und Nenner fallen gleich, das Produkt bleibt gleich. \quad c) $6 \cdot 0{,}4 \cdot 0{,}35 \cdot 0{,}25 = 0{,}21$ \quad d) „4 über 1“ · „5 über 2“ : „9 über 3“ $= \frac{40}{84} = \frac{10}{21}$}
